\section[Agent evaluation beyond workflow completion]{Agent evaluation beyond workflow completion: lessons from chemistry and software repair}
\label{sec:exec25-cow}

An agent that completes a tool workflow has not necessarily answered the scientific question. The distinction matters when a banking user asks whether a servicing action improves recoveries, an agent proposes a diagram, a worker returns an estimate, and a natural-language report describes a benefit. Each step can be locally plausible while the final answer addresses a different population, information set or monetary target. Two independently retained preprints provide complementary perspectives: Ramos, Collison and White organize scientific agents around observation, decision modules and tools; Zhou and colleagues examine vulnerability detection and repair, including data quality, evaluation and deployment limitations. These sources motivate a proposal design, not measured banking performance.\citep{exec25_cow_chemistry,exec25_cow_vulnerability}

The selected chemistry original is arXiv version3 submitted14November2024, with18November2024 also printed on its cover. The selected vulnerability survey is arXiv version3 of7October2024. Their journal identities are separately discoverable, but the retained original journal bytes were not obtained through the recorded routes and exact equivalence has not been established. Claims here refer to the inspected preprints. Their descriptions of earlier papers and packages are secondary accounts; this section does not inherit those underlying technical results without their own primary audit.\citep{exec25_cow_chemistry,exec25_cow_vulnerability}

\subsection{From business language to a scientific request}

The chemistry review defines a continuing observation/action loop, separates its trainable decision components from environment/tools, and groups memory, planning, reasoning and profiling as modules. The paper itself notes that agent nomenclature differs across frameworks. This taxonomy is a useful design vocabulary rather than a proof that one architecture is correct. For banking, we propose that every observation and action preserve an explicit scientific request
\[
 Q=(U,P,a,a_0,Y,h,\mathcal I_t,G,\mathcal A,C),
\]
where $U$ is the unit, $P$ the population, $a$ the contemplated action, $a_0$ its comparator, $Y$ the outcome, $h$ the horizon, $\mathcal I_t$ the available information, $G$ the candidate causal graph or graph family, $\mathcal A$ the assumptions, and $C$ the costs or constraints. The notation is our proposed contract. An LLM can elicit missing entries and propose alternatives; it cannot establish their truth by filling every field.\citep{exec25_cow_chemistry,exec25_cow_reference}

Suppose the business asks for net incremental value of a servicing action. In a stipulated example, the action adds expected gross recoveries USD70 and costs USD80, so its net increment is USD$-10$. Translating the question into gross recovery improvement gives a positive answer to a different decision. The correct response need not invent a missing cost: it can request a cost bound, return the threshold cost USD70, or refuse to recommend an action until the missing input is supplied. This semantic error can survive a syntactically valid graph, a numerically correct estimator and an articulate explanation. A public benchmark should therefore test conservation of the business target independently of diagram syntax and tool success.\citep{exec25_cow_reference}

Diagram construction should retain competing structures and reasons for each proposed edge. Source-derived claims, temporal ordering, domain assumptions and statistical suggestions need distinct provenance. Where observationally equivalent graphs imply different interventions, fluent graph completion is not identification. The LLM may ask which variables were recorded before assignment, whether policy changes affect observation, and whether linked customers interfere. Those questions become inspectable inputs to a causal tool. A structural restriction unsupported by the public fixture or a supplied domain premise should remain uncertain rather than becoming a high-confidence edge.\citep{exec25_cow_reference}

\subsection{Surrogate rewards and scientific validation}

The chemistry survey distinguishes predefined research assistance from hypothesis creation and discusses computational feedback in hypothesis refinement. It also notes unvalidated results in a cited generative example. Its own discussion of agent evaluation identifies unclear success criteria and proposes Bloom's taxonomy as one organizing approach. These are useful warnings about evaluation design; neither the taxonomy nor a predictor's score supplies a calibrated intervention-value guarantee.\citep{exec25_cow_chemistry}

For an exact counterexample, consider two proposed actions with proxy scores0.9and0.8 and true stipulated net utilities USD$-10$andUSD20. Selecting the higher score chooses the lower-utility action and incurs USD30regret relative to the better action. This example does not show that the survey's systems incurred that regret. It establishes why a proposed bank experiment must declare the mapping from proxy to target and measure its failure under controlled shifts. A simulator, reward model or natural-language judge can be a useful oracle only for the properties justified by its assumptions and validation.\citep{exec25_cow_reference}

Adding judges is not automatically independent evidence. Let each judge have marginal error $q=1/10$. Three independent judges have majority error
\[
 3q^2(1-q)+q^3=\frac{7}{250}=0.028.
\]
If all three copy the same mistaken premise, their common error probability remains0.1. Both constructions have the same marginal accuracy. A proposal should therefore qualify the dependence, shared context and training overlap of its evaluation channels, and separate a model critique from an externally specified causal oracle. These finite calculations are project examples, not measured error rates for any model.\citep{exec25_cow_reference}

The chemistry review also cautions that sanity-checking tools do not make an answer hallucination-proof. Its proposed narrative comparing lengthy laboratory exploration with a one-week AI-assisted route is explicitly a scenario, not a controlled productivity estimate. We should apply the same discipline to the banking investment case: separate an illustrative workflow from observed time, scientific acceptance and actual cost. Public evidence can justify an expensive prototype or bank experiment without pretending that its return has already been measured.\citep{exec25_cow_chemistry,exec25_cow_reference}

\subsection{What software-repair evidence contributes}

Zhou and colleagues' preprint reviews58primary studies within its stated search scope. Its synthesis identifies heuristic vulnerability labels, repair datasets lacking tests, limited input granularity and offline evaluation/deployment gaps. Those are findings of that dated review, not current universal absence claims. Its threats-to-validity discussion acknowledges terminology and search-selection limits. We use the review as primary evidence of its own synthesis and proposed roadmap; individual model-performance results remain dependent on their original studies.\citep{exec25_cow_vulnerability}

The connection to causal software is direct at the level of validation logic. A generated adapter may compile and pass a sampled test while reversing an intervention flag, dropping failed requests or joining future information. For a finite example, let the required function be $f(x)=x$ on $\{0,1\}$. A constant-zero patch passes a test at0and fails the specification at1. The number of passing tests cannot be interpreted without their coverage and oracle. For a separate label example, predictions can agree with all four heuristic reference labels and disagree with all four stipulated true labels. These are exact constructions, not allegations about a particular benchmark. They motivate independent scientific oracles and adverse cases for generated code.\citep{exec25_cow_reference}

The vulnerability survey's discussion of training/evaluation overlap also matters to a public causal-agent study. A proposal should bind task provenance, source versions, temporal or structural splits, and known exposure limits. Public accessibility makes a benchmark useful for reproducibility but does not establish that a model has never encountered it. A new natural-language paraphrase need not create a new causal problem. Splitting by graph family, mechanism, decision target or source document can address different leakage channels, but each split requires a declared generalization question. This is an experiment design to test, not a claim of contamination-free evaluation.\citep{exec25_cow_vulnerability,exec25_cow_reference}

\subsection{A staged public evaluation contract}

The first stage evaluates translation and diagrams on finite public/generated worlds with known units, clocks, assignments and outcomes. It scores the scientific target, allowed assumptions, graph family, identification or identified-set answer, and correct refusal separately. The second stage evaluates tool and code execution: typed requests, immutable inputs, output schemas, numerical tolerances, fault injection and complete denominators. The third stage evaluates decisions and costs under the same target law, followed by explicit transfer restrictions for a later bank pilot. Progress in one stage does not automatically pass the others.\citep{exec25_cow_reference}

For example,100requested tasks can produce80completed tool runs and60scientifically accepted answers. Completion is80\%, accepted answers per request60\%, and acceptance conditional on completion75\%. With a stipulated total cost USD200, cost per completion isUSD2.50but cost per accepted answer isUSD$10/3$. Valid refusals and unresolved outcomes should be reported separately rather than discarded. A scientific scoring rule must specify whether an honest refusal is a success for the requested task. None of these denominators is a measured value for a deployed agent.\citep{exec25_cow_reference}

At the CLI/MCP boundary, a proposed request includes the versioned $Q$, evidence references, permitted worker operations, resource limits and logical attempt identity. A result retains the original request alongside the graph/assumption decisions, estimator or refusal, scientific diagnostics, code/artifact hashes and complete tool outcomes. Logs should support reconstruction of why a semantic choice changed; a narrative summary cannot replace the underlying bindings. Generated code runs only in an explicitly qualified worker scope, and a successful interface call grants no scientific acceptance. No repository integration is implemented by this design discussion.\citep{exec25_cow_reference}

Known repair routes include explicit semantic fixtures, independently derived causal oracles, versioned source custody, reproducible numerical tests and complete cost accounting. Unavailable journal originals, full underlying-study inspection and task-specific contamination evidence remain bounded obligations. Unknown bank assignment truth, unrestricted latent identification and universal transfer of agent accuracy cannot be settled by this public survey. Modular tools make the process inspectable; scientific acceptance still depends on the target, evidence and validation law.\citep{exec25_cow_reference}
